\section*{ЗАКЛЮЧЕНИЕ}
\phantomsection
\addcontentsline{toc}{section}{Заключение}

В результате выполнения курсовой работы решена задача численного анализа влияния дискретного микрорельефа на трибологические характеристики радиального подшипника скольжения. Были выполнены следующие задачи:

\begin{enumerate}
    \item Проведён обзор теоретических основ гидродинамической смазки и современных исследований в области лазерного текстурирования поверхностей трения. Установлено, что ключевым механизмом положительного влияния микрорельефа является микрогидродинамический эффект, создаваемый локальными клиновыми зазорами на кромках углублений.

    \item Освоена методика численного решения уравнения Рейнольдса для подшипника конечной длины методом конечных разностей с итерационной процедурой Гаусса-Зейделя и верхней релаксацией. Расчёт выполнен на сетке $500 \times 500$ узлов с критерием сходимости $\delta < 10^{-5}$.

    \item Получены поля давления и интегральные характеристики (нагрузочная способность $F$, коэффициент трения $\mu$, расход смазки $Q$) для гладкого подшипника и подшипника с эллипсоидальными углублениями (тип~1) в диапазоне эксцентриситетов $\varepsilon = 0{,}05 - 0{,}80$. Геометрия подшипника: $R = 35$~мм, $c = 50$~мкм, $L = 56$~мм.

    \item Выполнен сравнительный анализ гладкого и текстурированного подшипников. При характерном рабочем эксцентриситете $\varepsilon = 0{,}6$ зафиксированы следующие изменения:
    \begin{itemize}
        \item нагрузочная способность увеличилась на $10{,}82$\%: с $F = 5028{,}5$~Н до $F = 5572{,}7$~Н;
        \item коэффициент трения снизился на $11{,}78$\%: с $\mu = 0{,}01538$ до $\mu = 0{,}01356$;
        \item расход смазки возрос на $1{,}43$\%: с $Q = 0{,}2359$~л/с до $Q = 0{,}2392$~л/с.
    \end{itemize}
\end{enumerate}

Полученные результаты подтверждают, что нанесение эллипсоидальных углублений с параметрами $a = 2{,}41$~мм, $b = 2{,}214$~мм, $h_p = 10$~мкм ($h_p/c = 0{,}2$) приводит к существенному улучшению рабочих характеристик подшипника. Одновременное повышение нагрузочной способности и снижение коэффициента трения при незначительном росте расхода смазки свидетельствует о высокой эффективности данного типа микрорельефа.

Эллипсоидальный профиль углублений, благодаря плавному изменению глубины от центра к периферии, обеспечивает формирование устойчивого дополнительного давления без резких скачков и разрывов потока. Это делает его предпочтительным выбором для подшипников, работающих в стационарных режимах при умеренных и высоких эксцентриситетах.

На основании проведённого анализа можно рекомендовать применение эллипсоидального микрорельефа для подшипников скольжения нефтегазового оборудования, работающего в условиях жидкостного трения. Особенно перспективным является использование текстурирования в тяжелонагруженных узлах ($\varepsilon > 0{,}5$), где прирост нагрузочной способности и снижение трения выражены наиболее значительно.

